% ============================================================================
% Fichier  : partie4/P04_C13_forensique_reseau.tex : FORN (LPI342)
% Sources : RFC 3227 ; NIST SP 800-86 ; MITRE ATT&CK. Opération TSANAGA.
% Le bureau d'analyse travaille sur des journaux de flux exportés (pas de
% capture graphique).
% ============================================================================

\chapter{Forensique Réseau}
\label{chap:reseau}
\chaptermark{Forensique réseau}

\epigraph{%
    « Le disque dit ce qu'on a gardé~; le réseau dit ce qu'on a envoyé. »%
}{Principe du cours}

\niveauChapitre{Intermédiaire à expert}{Chapitre~\ref{chap:memoire}~; notions
de réseaux.}

\begin{boxobjectifs}
\begin{itemize}
    \item Identifier les sources de traces réseau.
    \item Repérer une exfiltration et un canal de pilotage.
    \item Lire des journaux de flux exportés.
    \item Analyser les échanges suspects (pilotage, dissimulation par DNS).
    \item Reconstituer une chronologie réseau.
\end{itemize}
\end{boxobjectifs}

\bigskip

% ============================================================================
\section*{Question directrice : le fichier clients a-t-il quitté le réseau de LiZbiZ, et vers où ?}
% ============================================================================

Au-delà de la copie sur clé, M.~NDONGO a-t-il aussi envoyé le fichier par le
réseau~? Les journaux de flux et de mandataire gardent la trace des envois.
Ce chapitre apprend à les lire.

\medskip
\textbf{Plan du chapitre.}
\begin{enumerate}
    \item Les sources réseau (section~\ref{sec:c13-sources}).
    \item Exfiltration et pilotage (section~\ref{sec:c13-exfil}).
    \item Les échanges suspects (section~\ref{sec:c13-suspects}).
    \item La chronologie réseau (section~\ref{sec:c13-chrono}).
\end{enumerate}

\separateur

% ============================================================================
\section{Les Sources Réseau}
\label{sec:c13-sources}
% ============================================================================

\begin{tableaucours}{L{3.6cm}L{8.6cm}}{%
    \tableauHeader{Source} &
    \tableauHeader{Ce qu'elle révèle}
}
Journaux de flux (pare-feu) & Qui parle à qui, quand, quel volume \\
Journaux du mandataire web & Destinations web, volumes, pièces envoyées \\
Journaux DNS & Noms résolus, dissimulation par DNS \\
Journaux de messagerie & Courriels, pièces jointes, destinataires \\
\end{tableaucours}
\vspace{-0.8em}
\begin{center}{\small\itshape Tableau~13.1. Sources de traces réseau}\end{center}

\begin{boxalert}
\textbf{Contrainte du cours.} Le bureau d'analyse de NDAP SUKULU n'a pas
d'outil de capture graphique.

\smallskip
\textbf{Règle~:} on travaille sur des \emph{journaux de flux exportés} (au
format tableur), suffisants pour établir les faits et parfaitement opposables.
La méthode prime sur l'outil.
\end{boxalert}

% ============================================================================
\section{Exfiltration et Pilotage}
\label{sec:c13-exfil}
% ============================================================================

Deux signatures réseau comptent~: l'\termecle{exfiltration} (un envoi sortant
anormal vers une destination externe) et, dans les affaires d'intrusion, le
\termecle{canal de pilotage} (petites connexions régulières commandant un code
malveillant). Dans TSANAGA, on cherche surtout l'envoi du fichier clients.

% ============================================================================
\section{Les Échanges Suspects}
\label{sec:c13-suspects}
% ============================================================================

Certains envois se dissimulent~: données glissées dans des requêtes DNS,
transferts vers un service d'hébergement personnel, courriels à une adresse
externe. L'analyste repère les destinations inhabituelles et les volumes qui
ne correspondent pas à l'usage normal.

% ============================================================================
\section{La Chronologie Réseau}
\label{sec:c13-chrono}
% ============================================================================

\begin{boxscenario}{Opération TSANAGA. L'envoi du 5 janvier}
\textbf{Constatation.} Les journaux du mandataire montrent, le 5~janvier, un
envoi de 18~Mo depuis le poste de M.~NDONGO vers un service d'hébergement
personnel, deux jours après la copie sur clé.

\medskip
\textbf{Analyse} (confiance élevée pour l'envoi, moyenne pour le contenu).
L'envoi est établi~; que son contenu soit le fichier clients est fortement
probable (taille, date), mais se recoupe avec le disque. On date le fait, on
qualifie la probabilité.
\end{boxscenario}

\begin{boxmetier}{Analyste forensique}
\textbf{Suivre la donnée.} L'analyste trace le chemin de la donnée~: copiée
sur une clé, envoyée par le réseau. Chaque canal laisse une trace~; les
recouper reconstitue l'histoire complète de l'exfiltration.
\end{boxmetier}

\separateur

\begin{boxresume}
\textbf{Ce qu'il faut retenir de ce chapitre~:}
\begin{itemize}
    \item \textbf{Flux, mandataire, DNS, messagerie}~: les sources réseau.
    \item \textbf{Exfiltration} (envoi sortant) et \textbf{pilotage} (petites
          connexions régulières).
    \item \textbf{Des journaux de flux exportés} suffisent.
    \item \textbf{Les échanges se dissimulent}~: destinations inhabituelles.
    \item \textbf{On reconstitue} la chronologie et on recoupe.
\end{itemize}
\end{boxresume}

\begin{boxinfo}
\textbf{Points de contrôle de l'élève.} Je sais~: citer les sources réseau~;
reconnaître exfiltration et pilotage~; lire des flux exportés~; repérer un
échange suspect~; dater un envoi.
\end{boxinfo}

\bigskip

\begin{boxexo}[title={\ding{45}\quad Exercice~13.1. Sources \niveaubase}]
Quelle source révèle les pièces jointes envoyées par courriel~?
\end{boxexo}

\begin{boxexo}[title={\ding{45}\quad Exercice~13.2. Exfiltration \niveaubase}]
Quelle signature réseau trahit un envoi sortant anormal~?
\end{boxexo}

\begin{boxexo}[title={\ding{45}\quad Exercice~13.3. Sans capture \niveaubase}]
Pourquoi des journaux de flux exportés suffisent-ils~?
\end{boxexo}

\begin{boxexo}[title={\ding{45}\quad Exercice~13.4. Dissimulation \niveaumoyen}]
Comment des données peuvent-elles se dissimuler dans des requêtes DNS~?
\end{boxexo}

\begin{boxexo}[title={\ding{45}\quad Exercice~13.5. TSANAGA \niveaumoyen}]
Qu'établissent les journaux du mandataire le 5~janvier~?
\end{boxexo}

\begin{boxexo}[title={\ding{45}\quad Exercice~13.6. Contenu \niveaumoyen}]
Pourquoi l'envoi est-il « établi » mais son contenu « probable »~?
\end{boxexo}

\begin{boxexo}[title={\ding{45}\quad Exercice~13.7. Recouper \niveaumoyen}]
Comment recouper l'envoi réseau avec le disque~?
\end{boxexo}

\begin{boxexo}[title={\ding{45}\quad Exercice~13.8. Destination \niveaumoyen}]
Comment repère-t-on une destination inhabituelle dans les flux~?
\end{boxexo}

\begin{boxexo}[title={\ding{45}\quad Exercice~13.9. Chronologie \niveauexpert}]
Dressez la chronologie réseau de l'exfiltration dans TSANAGA.
\end{boxexo}

\begin{boxexo}[title={\ding{45}\quad Exercice~13.10. Chemin de la donnée \niveauexpert}]
Reconstituez le chemin complet du fichier clients (clé et réseau).
\end{boxexo}

\bigskip

\begin{boxinfo}
\textbf{Pour aller plus loin.}
\begin{itemize}
    \item RFC~3227~\cite{rfc3227}~; NIST SP~800-86~\cite{nist80086}~;
          MITRE ATT\&CK~\cite{mitreattack}.
\end{itemize}
\end{boxinfo}

\begin{boxinfo}
\textbf{Transition.} L'enquête se déplace vers la poche du suspect~: le
Chapitre~\ref{chap:mobile} traite du forensique mobile.
\end{boxinfo}
